\section{Modern Chatnotes: the four-label coefficient algebra}
\label{sec:chatnotes-coefficient-algebra}

This section is an editorial reconstruction of raw Chatnotes, not an
attribution to the authors of the external literature.  The controlling local
witness is the raw export \emph{Z\_n Symmetries of Collatz}, physical
lines~5712--6220.  The external source actually cited in that interval is
Aizawa--Isaac--Segar's construction of
$\mathbb Z_2\times\mathbb Z_2$ color superalgebras inside universal
enveloping algebras \cite{AizawaIsaacSegar2019}.  The elementary coefficient
results below are proved directly.  The cited paper controls only the color
definition and its own enveloping-algebra example.

\subsection{Four affine labels and the state-dependent evaluator}

Put
\[
 V=\mathbb F_2^2,
 \qquad
 F_{a,b}(X)=\frac{(1+2a)X+b}{2}
 \quad ((a,b)\in V).
\]
Thus the four labels give
\[
 F_{00}(X)=\frac X2,
 \quad F_{01}(X)=\frac{X+1}{2},
 \quad F_{10}(X)=\frac{3X}{2},
 \quad F_{11}(X)=\frac{3X+1}{2}.
\]
Let $\operatorname{Aff}_{\mathbb Q}^{\times}$ denote the group of affine
maps $X\mapsto mX+c$ with $m\in\mathbb Q^\times$, represented by
coordinates $(m,c)$ and composition
\[
 (m,c)(m',c')=(mm',mc'+c).
\]

\begin{proposition}[exact type of the four-label map]
\label{prop:chatnotes-affine-label-map}
The map of sets
\[
 \lambda:V\longrightarrow \operatorname{Aff}_{\mathbb Q}^{\times},
 \qquad
 \lambda(a,b)=\left(\frac{1+2a}{2},\frac b2\right),
\]
is injective.  It is not a homomorphism from the additive group $V$ to the
affine group under composition.  For every $n\in\mathbb Z$,
\[
 F_{a,b}(n)\in\mathbb Z
 \quad\Longleftrightarrow\quad
 n\equiv b\pmod 2.
\]
\end{proposition}

\begin{proof}
Equality of two affine maps gives equality of their intercepts, hence
$b=b'$, and equality of their slopes, hence $a=a'$.  This proves
injectivity.  The identity of the affine group is $(1,0)$, whereas the zero
of $V$ is sent to $(1/2,0)$.  A group homomorphism must preserve the identity,
so $\lambda$ is not one.  Finally, $1+2a$ is odd, and therefore
\[
 (1+2a)n+b\equiv n+b\pmod2.
\]
Division by two is integral exactly when this numerator is even, equivalently
$n\equiv b\pmod2$.
\end{proof}

Define
\[
 \partial:V\longrightarrow\mathbb F_2,
 \qquad \partial(a,b)=a+b,
 \qquad
 \Delta=\ker\partial=\{00,11\}.
\]
For $n\in\mathbb Z$, let $\overline n$ be its class in $\mathbb F_2$ and set
\[
 \iota(n)=(\overline n,\overline n)\in\Delta.
\]
The map $\iota:\mathbb Z\to\Delta$ is a surjective group homomorphism with
kernel $2\mathbb Z$ and fibres
\[
 \iota^{-1}(00)=2\mathbb Z,
 \qquad
 \iota^{-1}(11)=2\mathbb Z+1.
\]

\begin{proposition}[typed diagonal evaluation]
\label{prop:chatnotes-diagonal-evaluator}
Let
\[
 D=\{(n,v)\in\mathbb Z\times\Delta:v=\iota(n)\},
 \qquad
 \eta:\mathbb Z\longrightarrow D,
 \quad \eta(n)=(n,\iota(n)).
\]
Then $\eta$ is a bijection with inverse the first projection.  The formula
\[
 E:D\longrightarrow\mathbb Z,
 \qquad E(n,v)=F_v(n),
\]
is well-defined, and the shortened Collatz map is exactly
\[
 H=E\circ\eta,
 \qquad
 H(n)=
 \begin{cases}
 n/2,&n\equiv0\pmod2,\\
 (3n+1)/2,&n\equiv1\pmod2.
 \end{cases}
\]
\end{proposition}

\begin{proof}
The definition of $D$ makes the two inverse identities for $\eta$ and the
first projection immediate.  If $(n,v)\in D$, then
$v=(\overline n,\overline n)$, so Proposition~\ref{prop:chatnotes-affine-label-map}
proves $F_v(n)\in\mathbb Z$.  Substitution of $v=00$ on the even fibre and
$v=11$ on the odd fibre gives the displayed formula for $H$.
\end{proof}

The last proposition is a coefficient-label encoding with an exact evaluator.
It does not make $\lambda$ an action, and it does not turn iteration of $H$
into addition in $V$.

\subsection{Function algebra and group algebra without conflation}

There are two four-dimensional algebras in the raw passage.  Write
\[
 A=\operatorname{Map}(V,\mathbb C)
\]
with pointwise multiplication, and let $\delta_u$ be the point mass at
$u\in V$.  Then
\[
 \delta_u\delta_w=\mathbf 1_{u=w}\delta_u,
 \qquad
 1_A=\sum_{u\in V}\delta_u.
\]
Write
\[
 B=\mathbb C[V]
\]
for the group algebra, with group-element basis $g_v$.  Its multiplication is
\[
 g_vg_w=g_{v+w},
 \qquad 1_B=g_{00}.
\]
The point masses are therefore not the group-element basis.

Use the symmetric bilinear form
\[
 \langle u,v\rangle=u_1v_1+u_2v_2\in\mathbb F_2
\]
and define the Fourier idempotents
\[
 p_u=\frac14\sum_{v\in V}(-1)^{\langle u,v\rangle}g_v.
\]

\begin{theorem}[exact Fourier algebra isomorphism]
\label{thm:chatnotes-fourier-algebra-isomorphism}
The linear map
\[
 \mathcal F:A\longrightarrow B,
 \qquad \mathcal F(\delta_u)=p_u,
\]
is a unital algebra isomorphism.  Its inverse is fixed on the group basis by
\[
 \mathcal F^{-1}(g_v)
 =\sum_{u\in V}(-1)^{\langle u,v\rangle}\delta_u.
\]
Equivalently,
\[
 p_up_w=\mathbf 1_{u=w}p_u,
 \qquad
 \sum_{u\in V}p_u=g_{00},
 \qquad
 g_v=\sum_{u\in V}(-1)^{\langle u,v\rangle}p_u.
\]
\end{theorem}

\begin{proof}
For $z\in V$, the coefficient of $g_z$ in $p_up_w$ is
\[
 \frac1{16}\sum_{v\in V}
 (-1)^{\langle u,v\rangle+\langle w,z-v\rangle}
 =\frac{(-1)^{\langle w,z\rangle}}{16}
   \sum_{v\in V}(-1)^{\langle u-w,v\rangle}.
\]
Character orthogonality makes the last sum equal to $4$ when $u=w$ and
$0$ otherwise.  This proves $p_up_w=\mathbf 1_{u=w}p_u$.
Applying the same orthogonality relation to $\sum_up_u$ leaves only $g_{00}$.
It also gives the displayed expansion of every $g_v$ in the $p_u$ basis.
Thus the four $p_u$ form a basis, the proposed inverse identities hold, and
$\mathcal F$ preserves multiplication and the unit.
\end{proof}

This isomorphism loses no algebraic information, but it does not preserve the
natural $V$-grading
\[
 B=\bigoplus_{v\in V}\mathbb Cg_v,
 \qquad \deg g_v=v.
\]
Every $p_u$ has a nonzero component in all four homogeneous summands.
There is also a basis-independent obstruction to declaring the point masses
homogeneous in their labels.  If a nonzero idempotent $e$ in a direct
$V$-graded algebra had degree $u\ne00$, then $e^2$ would have degree
$u+u=00$, while $e^2=e$ would have degree $u$.  Directness of the grading
would force $e=0$, a contradiction.

\subsection{What the diagonal projector actually returns}

In $A\otimes_{\mathbb C}\mathbb C[X]$, define
\[
 \mathfrak F(X)=\sum_{v\in V}\delta_v\otimes F_v(X),
 \qquad
 P_\Delta=\delta_{00}+\delta_{11}.
\]
Pointwise multiplication gives
\begin{equation}
\label{eq:chatnotes-diagonal-projection}
 P_\Delta\mathfrak F(X)P_\Delta
 =\delta_{00}\otimes\frac X2
  +\delta_{11}\otimes\frac{3X+1}{2}.
\end{equation}
The right-hand side is a two-component labelled polynomial; it is not a
scalar map $\mathbb Z\to\mathbb Z$.  If
$\operatorname{ev}_v:A\to\mathbb C$ denotes evaluation at $v$ and
$\operatorname{ev}_{X=n}:\mathbb C[X]\to\mathbb C$ denotes substitution,
then the exact scalar recovery is
\begin{equation}
\label{eq:chatnotes-state-dependent-evaluation}
 H(n)=
 \bigl(\operatorname{ev}_{\iota(n)}\otimes
       \operatorname{ev}_{X=n}\bigr)
 \bigl(P_\Delta\mathfrak F(X)P_\Delta\bigr).
\end{equation}
The state-dependent evaluator in~\eqref{eq:chatnotes-state-dependent-evaluation}
cannot be omitted: projection removes the two off-diagonal labels but does
not choose between $00$ and $11$.

Under Theorem~\ref{thm:chatnotes-fourier-algebra-isomorphism}, the same
projector is
\begin{equation}
\label{eq:chatnotes-fourier-diagonal-projector}
 \mathcal F(P_\Delta)=p_{00}+p_{11}
 =\frac12(g_{00}+g_{11}).
\end{equation}
This is a proved change of basis, not an identification of the two basis
multiplications or gradings.

\subsection{The exact color-algebra consequence of the cited paper}

Aizawa--Isaac--Segar define a color superalgebra
$\mathfrak g=\bigoplus_{v\in V}\mathfrak g_v$ with bracket
\[
 \mathopen{[\![}X_v,Y_w\mathclose{]\!]}
 =X_vY_w-(-1)^{\langle v,w\rangle}Y_wX_v
\]
inside an enveloping algebra \cite{AizawaIsaacSegar2019}.  Starting from an
$\mathcal N=1$ superconformal Galilei algebra, they introduce
\[
 P_{nm}=\{P_n,P_m\},
 \qquad
 \Lambda_{nm}=\{P_n,X_m\},
 \qquad
 X_{nm}=[X_n,X_m]
\]
and assign degrees
\[
\begin{array}{c|l}
00&H,D,K,P_{nm},X_{nm},\\
01&P_n,\\
10&Q,S,\Lambda_{nm},\\
11&X_n.
\end{array}
\]
Their relation tables and Jacobi check establish closure for that specific
enveloping-algebra construction.

\begin{lemma}[totally isotropic diagonal restriction]
\label{lem:chatnotes-isotropic-diagonal}
Let $C=\bigoplus_{v\in V}C_v$ be an associative $V$-graded algebra.  Then
\[
 C_\Delta=C_{00}\oplus C_{11}
\]
is closed under multiplication.  On $C_\Delta$, the color commutator defined
by the displayed pairing is the ordinary commutator.
\end{lemma}

\begin{proof}
The set $\Delta$ is a subgroup, so $C_uC_v\subseteq C_{u+v}\subseteq
C_\Delta$ for $u,v\in\Delta$.  Moreover
$\langle00,00\rangle=\langle00,11\rangle=
\langle11,00\rangle=\langle11,11\rangle=0$ in $\mathbb F_2$.
Hence the color sign is $+1$ on every pair of diagonal homogeneous elements,
and the bracket is $XY-YX$.
\end{proof}

The lemma explains a valid formal feature of the repeated degree group.  It
does not supply a Collatz color superalgebra.  The raw passage defines no
graded algebra of Collatz states or paths, no bracket on such an algebra, no
degree-preserving bracket morphism to either algebra in the cited paper, and
no action of those algebras on Collatz states or observables.  The coefficient
algebra $A$ also cannot be given the asserted grading merely by declaring its
nonzero point idempotents to have their index degrees.

\subsection{The two involutions and the chosen orthogonal block}

Coordinate swap
\[
 s:V\longrightarrow V,
 \qquad s(a,b)=(b,a),
\]
fixes the diagonal pointwise.  In contrast, translation by $11$,
\[
 \tau(v)=v+11,
\]
exchanges $00$ and $11$.  Translation is an involution of the underlying
two-point diagonal set, but it is not a group automorphism because
$\tau(00)=11$.

On the real vector space
$W_\Delta=\mathbb R\delta_{00}\oplus\mathbb R\delta_{11}$, one may choose
\[
 S=\begin{pmatrix}0&1\\1&0\end{pmatrix},
 \qquad
 J=\begin{pmatrix}0&-1\\1&0\end{pmatrix}.
\]
Direct multiplication gives
\[
 S^2=I,
 \qquad J^2=-I,
 \qquad SJS^{-1}=-J,
 \qquad \det S=-1,
 \qquad \det J=1.
\]
Consequently $\{e^{\theta J}:\theta\in\mathbb R\}\cong SO(2)$ and adjoining
$S$ produces the standard semidirect product $SO(2)\rtimes\langle S\rangle
\cong O(2)$.

This matrix calculation is exact, but $J$ is additional structure.  Neither
$s$ nor the Collatz branch evaluator selects it; indeed $s$ fixes the two
diagonal labels rather than exchanging them.  No universal property,
dynamical action, lifting problem, or determinant anomaly is constructed by
the raw passage.  More precisely, the raw passage supplies neither a
homomorphism between the two ambient group systems nor a commutative diagram
relating the quotient map
\[
 q:\operatorname{Spin}(4)\times SU(2)
   \longrightarrow
   \bigl(\operatorname{Spin}(4)\times SU(2)\bigr)/\langle(-1,-1)\rangle
\]
to the coefficient homomorphism
\[
 \partial:V\longrightarrow\mathbb F_2,
 \qquad \partial(a,b)=a+b,
 \qquad \Delta=\ker\partial.
\]
The occurrence of order-two data in both constructions is therefore not, by
itself, a defined comparison arrow.  No identification or theorem transfer is
admitted without those arrows and their compatibility identities.

\begin{nonclaim}
The results of this section do not construct a
$\mathbb Z_2\times\mathbb Z_2$ action on Collatz dynamics, a Collatz color
superalgebra, an intrinsic $O(2)$ action or anomaly, a Spin--$SU(2)$ quotient
for Collatz, or a reformulation of the Collatz conjecture.  They establish the
exact coefficient encoding, the Fourier algebra isomorphism, the required
state-dependent evaluator, the isotropic-diagonal lemma, and the stated
obstructions at those scopes only.  The raw continuation turns to weighted
paths; its composition algebra must be read before the natural local programme
boundary is fixed.
\end{nonclaim}
